\section{Experiments}
\label{sec:experiments}

All numbers are read from \texttt{results/} and were produced by the full
pipeline on the complete dataset. Percentage points are abbreviated pp.

\subsection{Setup}
\label{sec:setup}

Hardware is an Apple M4 Pro with 48\,GB of memory and 12 threads. Seeds,
splits, and sample sizes are fixed in a single configuration file. The
development mode (a 5\,\% sample) runs the whole pipeline in minutes and is
used only for debugging; it is not reported.

\subsection{Classical tests and covariate balance}
\label{sec:classical}

On all 14M users, treated users visit more (4.854\,\% versus 3.820\,\%,
+1.034\,pp [1.006, 1.063]) and convert more (0.309\,\% versus 0.194\,\%,
+0.115\,pp [0.108, 0.122], a relative lift of +59.4\,\% [54.4, 64.7]). The
null $p_T=p_C$ is rejected for both outcomes ($p<10^{-170}$), and exact
bootstrap intervals (2,000 draws) match the Wald ones. These are the
\emph{naive} effects.

Under clean randomization, features should not predict treatment. They do,
weakly: a classifier two-sample test has AUC 0.509 [0.508, 0.510] (permutation
$p=0.001$) and a joint likelihood-ratio test gives $p=2\times10^{-53}$. The
imbalance is correlated with the outcomes, which is why the next section
adjusts for covariates.

\subsection{Average treatment effects}
\label{sec:ate}

Table~\ref{tab:ate} reports the ATE estimators. Every covariate-adjusted
estimator sits well below the difference in means, and they agree with each
other. The doubly robust estimate of the visit effect (+0.746\,pp) is 28\,\%
below the raw one; for conversion (+0.100\,pp) the reduction is 13\,\%.

\begin{table}[H]
\centering
\caption{ATE estimates in percentage points on all 14M users, with 95\,\%
confidence intervals. g-computation has no interval reported.}
\label{tab:ate}
\small
\begin{tabular}{l r r}
\toprule
\textbf{Estimator} & \textbf{Visit} & \textbf{Conversion} \\
\midrule
Difference in means & 1.034 [1.006, 1.063] & 0.1152 [0.1085, 0.1219] \\
Lin regression adjustment & 0.752 [0.726, 0.777] & 0.0990 [0.0923, 0.1056] \\
IPW, H{\'a}jek, LightGBM propensity & 0.761 [0.731, 0.791] & 0.1007 [0.0935, 0.1079] \\
\textbf{AIPW (doubly robust)} & \textbf{0.746 [0.721, 0.771]} & \textbf{0.1000 [0.0932, 0.1067]} \\
g-computation (S-learner) & 0.733 & 0.0977 \\
\bottomrule
\end{tabular}
\end{table}

\paragraph{Observational demonstration.} To check that adjustment does what
it claims, we make a confounded dataset out of the experiment: users are
dropped with a probability that depends on a feature and on treatment, over
16 replicates. The naive difference then reads 2.19\,pp for visits, with 0\,\%
interval coverage of the full-data AIPW answer. AIPW recovers 0.749\,pp
(coverage 88\,\%) and Lin regression 0.748\,pp.

\subsection{Assignment versus exposure}
\label{sec:exposure}

Only 3.6\,\% of ad-eligible users were shown an ad. The naive exposed-versus-control
contrast (+37.6\,pp on visits) is confounded by user activity: active users
are more likely to be exposed. Using randomized assignment as an instrument,
the covariate-adjusted effect of actually seeing an ad is +21.2\,pp
[20.5, 21.9] on visits and +2.84\,pp [2.65, 3.03] on conversions among
exposed users.

\begin{figure}[H]
\centering
\includegraphics[width=0.78\linewidth]{causal_exposure_iv.png}
\caption{Effect of actually seeing an ad, using assignment as an instrument,
compared with naive exposure contrasts.}
\label{fig:exposure}
\end{figure}

\subsection{Predictive models}
\label{sec:predictive-results}

\begin{table}[H]
\centering
\caption{Predictive models on the test split (2.8M users, base rate 0.29\,\%).
Brackets are 95\,\% bootstrap intervals.}
\label{tab:predict}
\small
\setlength{\tabcolsep}{5pt}
\resizebox{\linewidth}{!}{%
\begin{tabular}{l c c c c c}
\toprule
\textbf{Model} & \textbf{ROC-AUC} & \textbf{PR-AUC} & \textbf{Log loss} & \textbf{Brier} & \textbf{Pred./base} \\
\midrule
Logistic regression & 0.956 [0.953, 0.958] & 0.214 [0.205, 0.224] & 0.01202 & 0.00253 & 1.00 \\
Random forest & 0.953 [0.951, 0.956] & 0.227 [0.217, 0.236] & 0.01204 & 0.00250 & 0.99 \\
\textbf{LightGBM} & \textbf{0.958 [0.956, 0.960]} & \textbf{0.226 [0.217, 0.236]} & \textbf{0.01181} & 0.00251 & 1.00 \\
LightGBM, \texttt{scale\_pos\_weight}=342 & 0.873 & 0.044 & 1.332 & 0.086 & 51.1 \\
\bottomrule
\end{tabular}}
\end{table}

LightGBM beats logistic regression by a small but significant margin (paired
bootstrap: ROC-AUC +0.0022 [0.0017, 0.0027], PR-AUC +0.012 [0.008, 0.016])
and ties the random forest on PR-AUC. The unweighted models are already
calibrated: Platt and isotonic calibration change log loss by at most 0.0001.
Class weighting hurts ranking as well as calibration. Flagging the top 1\,\%
of users catches 54\,\% of converters at 15.7\,\% precision, 54 times the base
rate. Fourteen percent of test rows share a feature vector with a training
row, but they hold only 32 of 8,155 test conversions, and excluding them
leaves the metrics unchanged.

\begin{figure}[H]
\centering
\includegraphics[width=0.85\linewidth]{predict_roc_pr.png}
\caption{ROC and precision-recall curves. Precision-recall is the honest view
of a 0.3\,\% event; ROC-AUC looks strong because negatives dominate.}
\label{fig:rocpr}
\end{figure}

\subsection{Uplift models}
\label{sec:uplift}

Table~\ref{tab:qini} gives normalized Qini coefficients. On visits, the DR-
and S-learners beat the response model (+0.0052 [0.0022, 0.0082] and +0.0048
[0.0011, 0.0082]). On conversions, every learner ties or trails it, and the
causal forest is significantly worse ($-0.0079$ [$-0.0132$, $-0.0030$]).

\begin{table}[H]
\centering
\caption{Normalized Qini coefficient on the test split,
propensity-weighted, 95\,\% bootstrap intervals.}
\label{tab:qini}
\small
\begin{tabular}{l c c}
\toprule
\textbf{Scorer} & \textbf{Visit} & \textbf{Conversion} \\
\midrule
S-learner & 0.0777 [0.0723, 0.0836] & 0.1737 [0.1441, 0.2082] \\
T-learner & 0.0713 [0.0659, 0.0773] & 0.1339 [0.1040, 0.1640] \\
X-learner & 0.0773 [0.0718, 0.0842] & 0.1733 [0.1430, 0.2077] \\
DR-learner & 0.0781 [0.0721, 0.0846] & 0.1694 [0.1360, 0.2032] \\
Causal forest & 0.0740 [0.0689, 0.0803] & 0.1736 [0.1446, 0.2094] \\
Transformed outcome & 0.0745 [0.0685, 0.0803] & 0.1724 [0.1429, 0.2036] \\
Response model $P(Y\mid X)$ & 0.0729 [0.0671, 0.0793] & \textbf{0.1815 [0.1513, 0.2151]} \\
Random & $-0.0006$ & 0.0048 \\
\bottomrule
\end{tabular}
\end{table}

The BLP heterogeneity test finds real heterogeneity for every learner
($p\le2\times10^{-8}$), and the S-learner's visit slope is 1.02 [0.95, 1.09],
meaning well-calibrated effect scores. Negative predictions are noise: users
predicted to be \emph{hurt} by the ad (23\,\% of users under the T-learner)
have a positive observed effect (+0.13\,pp [0.003, 0.26] on visits).

\begin{figure}[H]
\centering
\includegraphics[width=0.85\linewidth]{cate_qini_curves.png}
\caption{Qini curves on the 2.8M-user test split.}
\label{fig:qini}
\end{figure}

\subsection{The central experiment: uplift versus response targeting}
\label{sec:central}

Table~\ref{tab:diff} gives incremental outcomes of uplift targeting minus
response targeting at matched budgets. For visits, uplift targeting at a
5\,\% budget finds 13,624 [12,612, 14,555] incremental visits against 8,542
[7,418, 9,674] for the response model, a gain of +5,083 [3,960, 6,098], about
60\,\%. Of the visits in the response model's top 5\,\%, 88\,\% would have
happened without the ad; for the uplift model's top 5\,\% the figure is 73\,\%.
The advantage persists at 10\,\% and is gone by 20--30\,\%. For conversions,
uplift is never significantly ahead.

\begin{table}[H]
\centering
\caption{Incremental outcomes on the test split, \textbf{uplift minus response
targeting} (95\,\% intervals). Visit uses the X-learner, conversion the
causal forest.}
\label{tab:diff}
\small
\begin{tabular}{l r r}
\toprule
\textbf{Budget} & \textbf{Visit} & \textbf{Conversion} \\
\midrule
Top 5\,\%  & \textbf{+5,083 [3,960, 6,098]} & $-42$ [$-151$, 65] \\
Top 10\,\% & \textbf{+2,862 [1,953, 3,840]} & $-76$ [$-162$, 24] \\
Top 20\,\% & +545 [$-303$, 1,335] & $-71$ [$-127$, 5] \\
Top 30\,\% & $-258$ [$-1{,}040$, 559] & $-53$ [$-104$, 9] \\
\bottomrule
\end{tabular}
\end{table}

\begin{figure}[H]
\centering
\includegraphics[width=0.85\linewidth]{cate_response_vs_uplift.png}
\caption{Uplift versus conversion-probability targeting, by budget.}
\label{fig:central}
\end{figure}

The share of conversions that would have happened anyway is flat, about
67\,\%, across score groups. The ad lifts conversion roughly in proportion to
baseline propensity, so the likeliest converters are also among the most
incremental.

\subsection{Targeting policies}
\label{sec:policies}

Table~\ref{tab:policy} compares policies on incremental conversions per 100K
users. Treating the top 10\,\% by conversion probability yields 89.5
[73.1, 104.3], against 100.4 [84.6, 114.7] for treating everyone: about 89\,\%
of the lift at a tenth of the treatment volume.

\begin{table}[H]
\centering
\caption{Incremental conversions per 100K users, test split,
propensity-weighted.}
\label{tab:policy}
\small
\resizebox{\linewidth}{!}{%
\begin{tabular}{l r r r r}
\toprule
\textbf{Policy} & \textbf{5\,\%} & \textbf{10\,\%} & \textbf{20\,\%} & \textbf{Treat all} \\
\midrule
Random & 5.9 [2.2, 9.5] & 12.3 [7.2, 16.8] & 22.9 [16.1, 29.5] & 100.4 [84.6, 114.7] \\
Highest conversion probability & 81.2 [67.2, 95.1] & 89.5 [73.1, 104.3] & 95.9 [81.1, 110.0] & \\
Highest estimated uplift & 79.7 [66.5, 93.2] & 86.8 [71.5, 101.2] & 93.3 [78.9, 107.2] & \\
Highest expected value & 79.7 & 86.8 & 92.5 (treats 16.7\,\%) & \\
\bottomrule
\end{tabular}}
\end{table}

With the budget chosen on validation, the economics favor targeting
(Table~\ref{tab:value}). At a cost-to-value ratio of 0.001, blanket treatment
nets about \$22 per 100K users, while response targeting at an 11\,\% budget
nets \$4,051 [3,330, 4,777].

\begin{table}[H]
\centering
\caption{Net value per 100K users, budget chosen on validation.}
\label{tab:value}
\small
\begin{tabular}{r l r r r}
\toprule
\textbf{Cost/value} & \textbf{Best policy} & \textbf{Budget} & \textbf{Net value} & \textbf{Treat all} \\
\midrule
0.0002 & Response & 31\,\% & \$4,628 & \$4,022 \\
0.0010 & Response & 11\,\% & \$4,051 [3,330, 4,777] & \$22 [$-790$, 758] \\
0.0020 & Response & 6\,\%  & \$3,634 & $-\$4{,}978$ \\
0.0050 & Response & 3\,\%  & \$2,775 & $-\$19{,}978$ \\
\bottomrule
\end{tabular}
\end{table}

\begin{figure}[H]
\centering
\includegraphics[width=0.85\linewidth]{targeting_incremental_vs_budget.png}
\caption{Incremental outcomes by treatment budget for each policy.}
\label{fig:targeting}
\end{figure}

\subsection{Robustness}
\label{sec:robustness}

Table~\ref{tab:robust} checks that the central result is not a seed or split
artifact. The visit advantage is stable (+2,893, SD 144, always positive over
three seeds) while the conversion deficit is small and consistently negative.
Conversion effect scores are much less stable than visit scores (rank
correlation between seeds 0.49 versus 0.98), as expected from a 0.3\,\% event.

\begin{table}[H]
\centering
\caption{Robustness of the uplift-versus-response comparison.}
\label{tab:robust}
\small
\resizebox{\linewidth}{!}{%
\begin{tabular}{l r r}
\toprule
 & \textbf{Visit (X-learner)} & \textbf{Conv. (causal forest)} \\
\midrule
Uplift $-$ response at 10\,\%, over 3 seeds & +2,893 (SD 144), always $>0$ & $-78$ (SD 14), always $<0$ \\
$\hat\tau$ rank correlation between seeds & 0.98 & 0.49 \\
$\hat\tau$ rank correlation between disjoint halves & 0.82 & 0.40 \\
Top-10\,\% set overlap between halves (Jaccard) & 0.70 & 0.68 \\
\bottomrule
\end{tabular}}
\end{table}

The X-learner's visit Qini passes the response model's only with at least 3M
training rows (Figure~\ref{fig:learning}), echoing the dataset authors'
finding that uplift needs very large samples. The pooled conversion model
ranks equally well in each randomized arm (ROC-AUC 0.959 control, 0.958
treated) but, never having seen treatment, over-predicts conversion within
control by 38\,\% and under-predicts within the treated arm by 4\,\%.

\begin{figure}[H]
\centering
\includegraphics[width=0.8\linewidth]{robustness_learning_curve.png}
\caption{Visit learning curve: how much data uplift ranking needs before it
overtakes the response model.}
\label{fig:learning}
\end{figure}

\subsection{Scaling}
\label{sec:scaling}

On all 8.39M training rows, logistic regression fits in 7.5\,s and peaks at
18.4\,GB, while LightGBM (300 trees) fits in 18.5\,s and peaks at 7.9\,GB,
with test PR-AUC 0.214 versus 0.224. LightGBM overtakes logistic regression on PR-AUC only above about
3M training rows. The full pipeline takes about 1.5--2 hours on this machine.
